% Clase 10 — las dos construcciones de las condiciones de borde, lado a lado.
% Van juntas a proposito: son dos superficies distintas (una cerrada, una curva)
% sobre la MISMA interfaz, y lo que las distingue es que una controla la
% componente normal y la otra la tangencial. Separarlas pierde el contraste.
% Las construcciones se dibujan MAS GRANDES que la version anterior: cuando la
% capsula media 1,2 x 0,8 los rotulos S1/S2/lateral no entraban y se montaban
% sobre la interfaz. Con mas aire cada uno tiene su lugar.
% La prosa explicativa va en el \caption; adentro solo las dos conclusiones.
\begin{tikzpicture}[scale=1]

  % ---------- panel izquierdo: capsula gaussiana -> D_n ----------
  \begin{scope}[shift={(0,0)}]
    \fill[figblue!8] (-2.5,0) .. controls (-1.1,0.34) and (1.1,-0.34) .. (2.5,0)
      -- (2.5,2.1) -- (-2.5,2.1) -- cycle;
    \fill[figamber!12] (-2.5,0) .. controls (-1.1,0.34) and (1.1,-0.34) .. (2.5,0)
      -- (2.5,-2.1) -- (-2.5,-2.1) -- cycle;
    \draw[figline, line width=1.0pt]
      (-2.5,0) .. controls (-1.1,0.34) and (1.1,-0.34) .. (2.5,0);
    \node[figlbl, figblue] at (-1.75,1.72) {medio 1};
    \node[figlbl, figamber] at (-1.75,-1.72) {medio 2};

    % capsula, con aire suficiente para los rotulos
    \draw[figred, line width=1.0pt] (-0.85,-0.62) rectangle (0.85,0.55);
    \node[figlblaux, figred, fill=figblue!8, inner sep=1.5pt] at (-1.35,0.55) {$S_1$};
    \draw[figaux] (-1.05,0.55) -- (-0.6,0.55);
    \node[figlblaux, figred, fill=figamber!12, inner sep=1.5pt] at (-1.35,-0.62) {$S_2$};
    \draw[figaux] (-1.05,-0.62) -- (-0.6,-0.62);
    \node[figlblaux, figred] at (1.78,0.55) {lateral};
    \draw[figaux] (1.4,0.48) -- (0.9,0.1);

    \draw[figvecres] (-0.42,0.55) -- (-0.42,1.18);
    \node[figlbl, figred] at (-0.75,1.02) {$\hat n$};

    \node[figlbl, align=center] at (0,-2.75) {Gauss para diel\'ectricos};
    \node[figlbl, align=center] at (0,-3.4)
      {$D_{1n}-D_{2n}=\sigma_{\text{libre}}$};
  \end{scope}

  % ---------- panel derecho: curva cerrada -> E_t ----------
  \begin{scope}[shift={(6.8,0)}]
    \fill[figblue!8] (-2.5,0) .. controls (-1.1,0.34) and (1.1,-0.34) .. (2.5,0)
      -- (2.5,2.1) -- (-2.5,2.1) -- cycle;
    \fill[figamber!12] (-2.5,0) .. controls (-1.1,0.34) and (1.1,-0.34) .. (2.5,0)
      -- (2.5,-2.1) -- (-2.5,-2.1) -- cycle;
    \draw[figline, line width=1.0pt]
      (-2.5,0) .. controls (-1.1,0.34) and (1.1,-0.34) .. (2.5,0);
    \node[figlbl, figblue] at (-1.75,1.72) {medio 1};
    \node[figlbl, figamber] at (-1.75,-1.72) {medio 2};

    % rectangulo de altura pequena
    \draw[figvecres] (-0.85,0.26) -- (0.85,0.2);
    \draw[figvecres] (0.85,0.2) -- (0.85,-0.28);
    \draw[figvecres] (0.85,-0.28) -- (-0.85,-0.22);
    \draw[figvecres] (-0.85,-0.22) -- (-0.85,0.26);
    \node[figlbl, figred, fill=figblue!8, inner sep=1.5pt] at (0,0.72) {$C$};
    \node[figlblaux, figred] at (1.85,0.62) {alto $\to 0$};
    \draw[figaux] (1.42,0.52) -- (0.95,0.12);

    \draw[figvecaux] (-0.85,-0.85) -- (0.85,-0.85);
    \node[figlblaux, figgray] at (0,-1.18) {$\hat l$};

    \node[figlbl, align=center] at (0,-2.75) {$\oint_C \vec E\cdot d\vec l=0$};
    \node[figlbl, align=center] at (0,-3.4) {$E_{1t}=E_{2t}$};
  \end{scope}

\end{tikzpicture}
